% !PRINTED 699
\weberpage{699}
Wäre nun $(O, B') > m$, so könnte man in $B'$ ein durch $\mathfrak{M}$
teilbares Polygon $\mathfrak{M}\mathfrak{B}$ finden (\S\,192, 5.), und es wäre:
\[
\mathfrak{A}' \mathfrak{M} \mathfrak{B} = \mathfrak{A} \mathfrak{B} = \mathfrak{W}
\]
ein vollständiges Polygon erster Gattung, es wäre also $\mathfrak{A}$ selbst
von der ersten Gattung; gegen die Voraussetzung. Also ist
$(O, B') = (A', W) = m$ und
\weq{15}{(O, A) = a - p + 1.}

Also gilt auch hier der \gesperrt{Riemann-Rochsche Satz} in der
Form 3.

Besteht $A$ aus einem isolierten Polygon $\mathfrak{M}$, so ist $(O, A) = 1$,
also $a = p$, und daraus folgt:

6. \gesperrt{Ein isoliertes Polygon zweiter Gattung mu\ss\
ein $p$-Eck sein, und umgekehrt ist jedes $p$-Eck
zweiter Gattung ein isoliertes Polygon.}

7. \gesperrt{Geht ein Punkt $\mathfrak{P}$ im Teiler einer uneigentlichen
 Klasse zweiter Gattung $A$ auf, und ist $A = \mathfrak{P} A'$,
so mu\ss\ $A'$ von der ersten Gattung sein.}

Denn es ist nach dem Satz (6) und (15):
\[
\begin{aligned}
(O, A) &= a - p + 1, \\
(O, A') &= a - p + (A', W), \\
(O, A) &= (O, A'),
\end{aligned}
\]
also $(A', W) = 1$.

\begin{center}
\textbf{\S\,200. Differentiale zweiter und dritter Gattung.}
\end{center}
Wir können jetzt die Bedingung für ein Differential:
\weq{1}{dJ = \frac{\mathfrak{A}}{\mathfrak{B}},}
die wir in \S\,197, (6) mit Rücksicht auf eine Variable $z$ in der
Form ausgedrückt haben:
\[
\mathfrak{U}^2 \mathfrak{A} \sim \mathfrak{Z} \mathfrak{B},
\]
invariant darstellen, wenn wir die vollständigen Polygone erster
Gattung benutzen, wonach
\[
\mathfrak{U}^2 \mathfrak{W} \sim \mathfrak{Z}
\]
ist:
\[
\mathfrak{A} \sim \mathfrak{W} \mathfrak{B},
\]
oder, indem wir die Klassen einführen:
%CONT>>
